\documentclass[landscape,a5paper,10pt]{article}
\usepackage[spanish]{babel}
\usepackage{ucs
			,amsfonts
			,amssymb
			,acronym
			,verbatim
			,rotating
			,graphicx
			,enumerate
			,fancyhdr
            ,mathrsfs
            ,array
			}

\usepackage{tikz}
\usetikzlibrary{svg.path}

\pagestyle{fancy}
\addtolength{\textheight}{2.0in}
\addtolength{\textwidth}{1.5in}
\addtolength{\voffset}{-1.0in}
\addtolength{\hoffset}{-1.0in}

\fancyhf{}
\fancyhead[L]{{\bfseries Programaci\'on I Curso 2012. \hfill ``Parcialito''}}
%\fancyfoot[LE,RO]{\thepage}

\fancypagestyle{firstpage}{
\fancyhf{}

\fancyhead[L]{\includegraphics[width=1.75cm]{ucalp.jpg}  \\
\small{UCALP. Facultad de Ciencias Exactas e Ingenier\'ia} \\
{\bfseries Programaci\'on I Curso 2012. \hfill ``Parcialito''}}
%\fancyfoot[LE,RO]{\thepage}
}

\thispagestyle{firstpage}

\begin{document}

~\\ \\ \\ \\

\begin{enumerate}

%\item Calcular el termino en\'esimo de la serie de Fibonacci con la ayuda de un arreglo de enteros.\\
%Serie de Fibonacci: ($f[0] = 1, \, f[1] = 1, \, f[n] = f[n-1] + f[n-2]$).

\item Dados dos arreglos de n\'umeros enteros previamente ordenado de longitud \verb=N= (constante), cree un nuevo arreglo con los datos de los dos arreglos mezclados de manera que queden igualmente ordenados. Luego imprima por pantalla el arreglo resultante.

\item Dadas dos matrices \verb=A= y \verb=B= de \verb=N x N= elementos:
\begin{itemize}
  \item[a-] Calcular la suma de \verb=A= y \verb=B=. ($C_{ij} = A_{ij} + B_{ij}$)
  \item[b-] Calcular el producto de \verb=A= por \verb=B=. ($C_{ij} = \sum A_{ij} * B_{ji} $)
\end{itemize}

\item Hacer un programa que lea de teclado dos fechas (representar el tipo de dato Fecha como un registro (d\'ia, mes, a\~no)) e imprima por pantalla la distancia entre dos fechas (en d\'ias, en meses, y en fechas).
    
%\item Representar el tipo de dato \emph{Fecha} como un registro (\emph{d\'ia, mes, a\~no}) y escribir un programa que:
%\begin{itemize}
%  \item[a-] Ordene un conjunto de fechas.
%  \item[b-] Permita calcular la distancia entre dos fechas (en d\'ias, en meses, y en fechas).
%\end{itemize}

\item Crear un tipo de dato \emph{conjunto extendido} que permita permita guardar (de alg\'un modo) m\'as de un elemento del mismo tipo y escribir un programa que permita la carga de datos y los muestre.

\item Desarrollar un programa que convierta una cantidad entera de d\'ias al formato a\~no, mes y d\'ia.

\item
    \begin{itemize}
        \item [a-] De ejemplos de alguna estructura de datos homog\'enea y otro de una heterog\'enea.
        \item [b-] Defina el \verb=type= de un conjunto de reales.
        \item [c-] De un ejemplo de un tipo definido por el usuario.
    \end{itemize}
%\begin{tabular}{p{0.4\textwidth} p{0.4\textwidth}}
%    \begin{itemize}
%        \item[a-] \begin{verbatim} mx7 := 924;\end{verbatim}
%        \item[c-] \begin{verbatim} rango := Gamma;\end{verbatim}
%        \item[e-] \begin{verbatim} rango := respuesta;\end{verbatim}
%         \item[g-] \begin{verbatim} tar := mx9;\end{verbatim}
%         \item[i-] \begin{verbatim} mx9 := Z;\end{verbatim}
%    \end{itemize}
%&
%    \begin{itemize}
%       \item[b-] \begin{verbatim} mx9 := Pagina;\end{verbatim}
%       \item[d-] \begin{verbatim} Gamma := tar;\end{verbatim}
%        \item[f-] \begin{verbatim} mx7 := Gamma;\end{verbatim}
%         \item[h-] \begin{verbatim} simbolo := `t`;\end{verbatim}
%          \item[j-] \begin{verbatim} respuesta := VERDADERO;\end{verbatim}
%    \end{itemize}
%\\
%\end{tabular}

\end{enumerate}
\end{document}
